%% appC --- The mathematics you need, on four pages
%%
%% A reference, not a lesson: each tool is introduced properly in the chapter
%% named beside it, and this appendix is where you look it up again. The few
%% constants it quotes come from code/measure/appc.py.
%%
%% Twin: appendices/pl/appC-toolkit.tex.

\chapter{The mathematics you need, on four pages}
\label{app:toolkit}

\noindent
Every tool the book uses, in one place, with the chapter that introduces it.
If a symbol stops you in the middle of a chapter, it is here.

\section{Very large and very small numbers}
\label{sec:appc-numbers}

\noindent
A number like $\num{0.00000000015}$ is hard to read, so it is written
$\num{1.5} \times 10^{-10}$: the digits that matter, and how many places to
move the decimal point. $10^{3}$ is a thousand, $10^{6}$ a million, and a
negative power moves the other way, so $10^{-3}$ is a thousandth. When the
book says two numbers differ \enquote{in the tenth decimal place}, it means
by something like $10^{-10}$. The book rounds what it prints and says how
many figures it keeps: a result is only worth the digits the measurement
behind it can support (Chapter~\ref{ch:models}).

\section{Rules, steps and orbits}
\label{sec:appc-steps}

\noindent
A rule that turns the present into the next moment is written
\[ x_{n+1} = f(x_{n}), \]
read \enquote{the next $x$ is $f$ of this $x$}. The small $n$ below the line
counts the steps: $x_{0}$ is the start, $x_{1}$ the state after one step, and
so on. The whole list $x_{0}, x_{1}, x_{2}, \ldots$ is called an
\emph{orbit}. A \emph{fixed point} is a state the rule leaves where it is,
a solution of $x = f(x)$ (Chapter~\ref{ch:iteration}).

\section{Powers, growth and halving}
\label{sec:appc-powers}

\noindent
Multiplying by the same factor $a$ again and again gives powers:
$a \times a \times a = a^{3}$, and $n$ steps of multiplying by $a$ multiply by
$a^{n}$. If $a$ is bigger than $1$ the result grows, faster and faster; if it
is between $0$ and $1$ it shrinks towards zero. Doubling ten times multiplies
by $2^{10} = 1024$, a little over a thousand, which is worth remembering: ten
doublings make a thousand, twenty make a million, and thirty a billion
(Chapters~\ref{ch:iteration}, \ref{ch:feedback} and~\ref{ch:butterfly}).

\section{Logarithms}
\label{sec:appc-logs}

\noindent
A logarithm answers the question \enquote{how many times?}. The natural
logarithm, written $\ln$, turns multiplying into adding:
\[ \ln(a \times b) = \ln a + \ln b . \]
So if an error is multiplied by $2$ at every step, $\ln$ of the error grows
by $\ln 2 \approx \val{appc.ln2}$ at every step, a steady amount, and a
steady amount is something you can average and compare. The number whose
$\ln$ is $1$ is called $e \approx \val{appc.e}$, and $\ln 10 \approx
\val{appc.ln10}$.

The question \enquote{how many doublings multiply by $N$?} has the answer
$\ln N / \ln 2$. Multiplying by a thousand takes about \val{appc.doublings.thousand}
doublings, and one extra decimal digit of precision is worth about
\val{appc.doublings.per.digit} doublings. That single fact is behind the
forecasting horizon of Chapter~\ref{ch:lyapunov} and the book's whole
argument about prediction: precision is paid for in digits, and each digit
buys only a fixed number of steps.

\section{The slope of a curve}
\label{sec:appc-slope}

\noindent
How steep is a curve at a point? Move the input a tiny amount $h$ either
side, see how much the output changes, and divide:
\[ \text{slope of } f \text{ at } x \approx \frac{f(x+h) - f(x-h)}{2h} . \]
For $f(x) = x^{3}$ at $x = 2$ the exact slope is $12$, and with
$h = \val{appc.slope.h}$ the recipe gives \val{appc.slope.approx}. A slope
tells you how much an error in the input is magnified in the output: a slope
of $3$ triples it, a slope of $-\tfrac{1}{2}$ halves it and flips its sign.
The size of a number without its sign is written with bars, $\abs{-3} = 3$
(Chapter~\ref{ch:iteration}).

\section{Rates of change}
\label{sec:appc-rates}

\noindent
When time runs smoothly rather than in steps, a rule says how fast each
quantity is changing right now. That is written
\[ \frac{\dd x}{\dd t} = \text{(some formula)}, \]
read \enquote{the rate of change of $x$ with time}. You never need to solve
such a rule with algebra in this book: a computer follows it forward in
many small steps of time, and Chapter~\ref{ch:motion} shows how to do that
well.

\section{Sums and averages}
\label{sec:appc-sums}

\noindent
A long sum is written with a capital Greek sigma:
\[ \sum_{n=1}^{N} a_{n} = a_{1} + a_{2} + \cdots + a_{N} , \]
and the average of $N$ numbers is that sum divided by $N$. The Lyapunov
exponent of Chapter~\ref{ch:lyapunov} is an average of this kind: the
average of $\ln$ of how much each step stretches a small error.

\section{Remainders}
\label{sec:appc-mod}

\noindent
\enquote{$x \bmod 1$} means \enquote{throw away the whole-number part and keep
what is after the decimal point}, so $\num{2.75} \bmod 1 = \num{0.75}$.
Chapter~\ref{ch:butterfly} uses it to keep a number between $0$ and $1$.

\section{Complex numbers}
\label{sec:appc-complex}

\noindent
A complex number is a point in a flat plane, written $a + b\iu$: $a$ steps
across and $b$ steps up. Adding two of them adds the steps. Multiplying uses
ordinary algebra plus one new rule, $\iu \times \iu = -1$, so
\[ (1 + 2\iu)(3 + 4\iu) = 3 + 4\iu + 6\iu + 8\iu^{2} = -5 + 10\iu . \]
The distance of $a + b\iu$ from the centre of the plane is
$\sqrt{a^{2} + b^{2}}$ (Chapter~\ref{ch:mandelbrot}).

\section{The Greek letters}
\label{sec:appc-greek}

\noindent
The book uses a handful, always for the same thing:

\begin{center}
\begin{tabular}{lll}
\toprule
Letter & Name & What it is in this book \\
\midrule
$\lambda$ & lambda & the Lyapunov exponent (Chapter~\ref{ch:lyapunov}) \\
$\delta$ & delta & Feigenbaum's constant (Chapter~\ref{ch:bifurcations}) \\
$\sigma, \rho, \beta$ & sigma, rho, beta & Lorenz's three settings (Chapter~\ref{ch:lorenz}) \\
$\varepsilon$ & epsilon & a small error or a small nudge \\
$\Sigma$ & capital sigma & a sum \\
\bottomrule
\end{tabular}
\end{center}
